\section{Every genuine finite model FS family lifts to actual original ample sections}
\label{sec:native-ample-model-fs-realization}

\begin{theorem}[Whole original ample-line realization of every model datum]
\label{thm:native-ample-model-fs-realization}
\lean{EnvelopeProperty.NativeAmpleModelFSRealization.ofModelData,EnvelopeProperty.NativeAmpleModelFSRealization.ofModelData_eval,EnvelopeProperty.NativeAmpleModelFSRealization.ofModelData_isRational}
\leanok
\uses{thm:native-ample-model-section-block-map,def:native-line-whole-section-fs,thm:native-ample-model-fs-evaluation,def:gl-fs}
Assume throughout this statement that the original scheme $X$ is integral.
Fix the same genuine original positive-power embedding data with $r>0$.
For each nonnegative rational original class $a$, every genuine model FS datum
of class $a/r$ is realized by actual entire original tensor-power sections of
$L$, with the same original positive denominator and the same weights. Its
whole weighted evaluation on every actual original native point equals the
original model evaluation under the same analytic homeomorphism. Original
rational weights remain rational.
\end{theorem}
\begin{proof}
The model degree $d$ is nonnegative because its denominator is positive and
$a/r\ge0$. The original whole linear block equivalence lifts every model
section of degree $d$ to an actual original section of $L^{\otimes rd}$.
Its degree equation is $rd=ma$, so the same original denominator and weights
give genuine original line FS data. Take the finite image of all the original
model terms under this actual entire inverse equivalence. Every lifted section
retains its true original model norm as the original genuine model section
norm. Thus each logarithmic term, the entire finite supremum and the whole
weighted function agree. Rationality holds with the same original weights.
No whole-section lift or finite-family realization certificate is an input.
\end{proof}

\begin{theorem}[Exact finite rational FS comparison on the actual original ample ray]
\label{thm:native-ample-model-finite-fs-iff}
\lean{EnvelopeProperty.NativeAmpleModelFSRealization.ofModelData_noCommonZero_iff,EnvelopeProperty.NativeAmpleModelFSRealization.isFinFS_iff}
\leanok
\uses{thm:native-ample-model-fs-realization,thm:native-ample-model-finite-fs-transport,def:gl-fs}
Assume throughout this statement that the original scheme $X$ is integral.
For $a\ge0$, a function on the entire original native point space is an actual
finite rational FS function of the original line $L$ and class $a$ if and only
if its transport under the same homeomorphism is a genuine model finite FS
function of class $a/r$. The actual lifted finite family has no common zero
exactly when its original model family has no common zero, on both entire
point spaces.
\end{theorem}
\begin{proof}
The exact true norm identity for each actual lifted section gives pointwise
nonvanishing in both directions. Use the actual finite image and surjectivity
of the entire native homeomorphism to compare the no-common-zero conditions
on both whole spaces. Combine this with unchanged rationality and the proved
whole function identity. The forward implication is the preceding actual
section-power construction; the reverse implication uses this genuine whole
linear section realization. The statement concerns the actual fixed original
ample-line ray and does not assert independence of an arbitrary numerical
representative or the complete global psh or envelope theorem.
\end{proof}
